\documentclass[a4paper]{article}
\usepackage{luatexja}
\usepackage{amsmath}
\usepackage{amssymb}
\usepackage{amsfonts}
\usepackage{bm}
\usepackage{geometry}
\geometry{left=25mm,right=25mm,top=25mm,bottom=25mm}
\usepackage{hyperref}

\title{射影幾何代数における二重時空：双対性、零並進、モーター}

\author{https://github.com/hypernumbernet}
\date{\today}

\renewcommand{\refname}{参考文献}

\begin{document}

\maketitle

\begin{abstract}
二重時空理論は、重力を、各質量粒子が運ぶ二つのミンコフスキー枠のあいだのねじれ不整合として扱う。それらの枠は双四元数代数 \(\mathrm{Cl}(3,1)\) のなかにある。そこにまだないのは並進である。射影幾何代数 \(G(3,1,1)\) は、単一の零ベクトル \(e_{4}\) を付加することでその欠落を埋める。並進はブーストや回転と同じ資格のバイベクトルになり、\(G(3,1,1)\) の10次元バイベクトル空間は、次元の数え上げではなく交換子についての事実として、ポアンカレ代数 \(\mathfrak{so}(3,1)\ltimes\mathbb{R}^{3,1}\) である。

構成を組織するのは三つの主張である。双対性 \(X\mapsto Xi\) はローレンツセクター上の複素構造である。すべてのブーストは回転の双対であり、\(\mathfrak{so}(3,1)\) を \(\mathfrak{so}(4)\) から分ける符号は \(i^{2}=-1\) である。ねじれと並進を一つにする幾何学的対象は、単一の Banach 指数ではなく、ねじれローターと零トランスレータの順序積 \(M=RT\) である。座標チャート上では、外部シュヴァルツシルト幾何はそのモーターの動径ブーストである。パラメータ空間で不整合を測る有限角のスカラー \(J\) は、遠平行の場の密度 \(T\) ではない。
\end{abstract}

\section{はじめに}
\label{sec:introduction}

二重時空理論は、共有された連続体を拒否する。各質量粒子は一対のミンコフスキー枠を運ぶ。通常枠はわれわれが測定する枠であり、双対枠はそれに対して時間反転し、同じ間隔をもつ。重力はその対のねじれである~\cite{Hestenes2015,DoranLasenby2003}。対を保持する代数は、\(\mathrm{Cl}(3,1)\) に標準的に同型な16次元双四元数代数である。通常セクターのブーストと双対セクターの回転は六つのバイベクトルによって生成される。それらの不整合は単一の代数的対象であり、その対象は、曲率を原始的なものとして取ることなく、球対称質量の外部を回復するのに足りる。

六つのバイベクトルは十個ではない。ローレンツ代数 \(\mathfrak{so}(3,1)\) は向きを説明する。変位は説明しない。粒子を一つの事象から別の事象へ動かすには、代数は並進も生成しなければならない。幾何代数においてその生成元の自然な住処は零ベクトルである。平方が零の方向を四つのミンコフスキー生成元と外積すれば、それ自身平方が零である四つのバイベクトルが得られる。これは射影幾何代数 \(G(3,1,1)\) の標準的な規約であり~\cite{Gunn2017,Lengyel2024,DeKeninckDorst2019,GunnDeKeninck2020}、退化クリフォード代数における剛体運動の Selig の扱いに遡る~\cite{Selig2005}。\(\mathrm{Cl}(3,1)\) に一つの零ベクトル \(e_{4}\) を付加すると、二重時空は32次元の代数へ持ち上がり、そのバイベクトル空間は10次元である。それら十個は三つのブースト、三つの回転、四つの並進である。

この持ち上げは帳簿上の操作ではない。交換子の下で10次元のスパンは閉じ、並進は可換イデアルをなし、ローレンツセクターはそのイデアルにベクトル表現で作用する。結果は \(G(3,1,1)\) の内部で実現されたポアンカレ代数 \(\mathfrak{so}(3,1)\ltimes\mathbb{R}^{3,1}\) である~\cite{Weinberg1995,Tung1985}。親理論ですでに擬スカラー \(i\) による右乗法として存在する双対性は、ローレンツ因子上の複素構造になる。ブーストは双対化された回転であり、\(i^{2}=-1\) がローレンツをユークリッドから区別するブースト同士の符号を強制する。対照的に並進は次数2の外へ送られる。双対性は向きを組織する。変位は組織しない。

両セクターが手元にあれば、モーターは順序積 \(M=RT\) である。ねじれ因子 \(R\) は6次元不整合の指数であり、並進因子 \(T\) は一次である。零バイベクトルの平方が零だからである。ユニタリ性 \(M\widetilde{M}=1\) は、すべてのバイベクトルが反転奇数であることから従う。それゆえ並進はベクトルではなくバイベクトルでなければならない。ねじれと並進が可換でないとき、\(M\) はそれらの和の指数ではない。幾何学的対象は積である。

同じモーターを動径ブーストに制限すれば、対角シュヴァルツシルトコフレームが回復される。それがしないのは、パラメータ空間のスカラー \(J\) を遠平行ねじれ密度 \(T\) と同一視することである。一方は有限角であり、他方はコフレームの導関数である。両者を混同すれば、定数ブーストが曲がった幾何に潰れ、あるいは純並進が非零の場になる。

論文はその順序に従う。第\ref{sec:adjunction}節は \(e_{4}\) を付加し、十個の生成元を分類する。第\ref{sec:duality}節は、双対性がローレンツ代数を定めることを示す。第\ref{sec:motors}節はモーターと二次診断量 \(J^{(5)}\) を構成する。第\ref{sec:gravity}節はシュヴァルツシルト外部を動径ブーストとして読む。第\ref{sec:conclusions}節は、代数が決着させたこととさせていないことへ戻る。それらの節の代数的恒等式は Lean~4 伴プロジェクト \texttt{DstDiophantine} において機械検証されている。対応する主張は付録~\ref{app:lean} に記録する。

\section{零方向の付加}
\label{sec:adjunction}

親理論は \(\mathrm{Cl}(3,1)\) の四つのベクトル生成元を双四元数要素
\[
e_{0}=j,\qquad e_{1}=kI,\qquad e_{2}=kJ,\qquad e_{3}=kK
\]
と同一視し、ミンコフスキー符号 \((-,+,+,+)\) を満たす：
\[
e_{0}^{2}=-1,\qquad e_{1}^{2}=e_{2}^{2}=e_{3}^{2}=+1,\qquad \{e_{\mu},e_{\nu}\}=0\quad(\mu\neq\nu).
\]
射影幾何代数 \(G(3,1,1)\) は、さらなる生成元 \(e_{4}\) を
\begin{equation}
e_{4}^{2}=0,\qquad \{e_{4},e_{\mu}\}=0\quad(\mu=0,1,2,3)
\label{eq:e4-defs}
\end{equation}
によって付加して得られる。完全なクリフォード代数は \(2^{5}=32\) 次元である。もとの16次元双四元数代数は \(\{e_{0},e_{1},e_{2},e_{3}\}\) が生成する部分代数である。

親理論の双対写像は、双四元数的擬スカラー~\cite{ConwaySmith2003}
\[
i \;=\; e_{0}e_{1}e_{2}e_{3} \;=\; j\,(kI)(kJ)(kK)
\]
による右乗法である。\(e_{4}\) は四つのミンコフスキーベクトルの各々と反可換であるから、それらの積とは可換である：
\[
e_{4}\,i \;=\; e_{4}\,e_{0}e_{1}e_{2}e_{3} \;=\; (-1)^{4}\,e_{0}e_{1}e_{2}e_{3}\,e_{4} \;=\; i\,e_{4}.
\]
したがって双対性は以前と全く同様に \(\mathrm{Cl}(3,1)\) 上で働き、新たな零方向は無傷のまま残る。持ち上げは \(X\mapsto Xi\) の時間反転的性格を乱さない。

\(G(3,1,1)\) のバイベクトル空間は次元 \(\binom{5}{2}=10\) をもつ。それら十個の要素は、平方によって区別される三つの幾何学的な類に分かれる。

三つの\emph{双曲型}生成元は平方が \(+1\) であり、\(\mathfrak{so}(3,1)\) のブーストを張る：
\[
iI \;=\; e_{0}e_{1},\qquad iJ \;=\; e_{0}e_{2},\qquad iK \;=\; e_{0}e_{3}.
\]
三つの\emph{巡回}生成元は平方が \(-1\) であり、同一のローレンツコピーの回転を張る：
\[
I \;=\; e_{3}e_{2},\qquad J \;=\; e_{1}e_{3},\qquad K \;=\; e_{2}e_{1}.
\]
四つの\emph{零}生成元は \(e_{4}\) とミンコフスキー基底との外積である、
\begin{align*}
N_{0} &= e_{4}\wedge e_{0} = e_{4}\,j, &
N_{1} &= e_{4}\wedge e_{1} = e_{4}\,(kI),\\
N_{2} &= e_{4}\wedge e_{2} = e_{4}\,(kJ), &
N_{3} &= e_{4}\wedge e_{3} = e_{4}\,(kK).
\end{align*}
\eqref{eq:e4-defs} から、任意の対 \(\mu,\nu\in\{0,1,2,3\}\) に対して
\begin{equation}
N_{\mu}N_{\nu} \;=\; (e_{4}e_{\mu})(e_{4}e_{\nu}) \;=\; -\,e_{4}^{2}\,e_{\mu}e_{\nu} \;=\; 0
\label{eq:null-strong-vanishing}
\end{equation}
が従う。各 \(N_{\mu}\) は平方零であり、混合積もすべて消える。四つの零バイベクトルは可換である。それらが並進になる。

光的並進は生成元への制約ではない。一般のミンコフスキーベクトル \(L=ct\,e_{0}+x\,e_{1}+y\,e_{2}+z\,e_{3}\) は \(L^{2}=0\) のとき零である。\(L\) に沿う並進はバイベクトル \(e_{4}\wedge L\) であり、四つの \(N_{\mu}\) の線形結合である。それは \(L\) 自身が零であるか否かに依らず、\eqref{eq:null-strong-vanishing} によって平方零である。光的性格は四つのパラメータ \(\lambda^{\mu}=(ct,x,y,z)\) の上に、錐 \(\eta_{\mu\nu}\lambda^{\mu}\lambda^{\nu}=0\) として住む。並進セクターは四次元のままである。

\begin{table}[htbp]
\centering
\caption{\(G(3,1,1)\) の10個のバイベクトル生成元。}
\label{tab:generators}
\begin{tabular}{llcl}
\hline
類     & 生成元 & 個数 & 平方 \\
\hline
双曲型 & \(B_{1}^{+}=iI,\; B_{2}^{+}=iJ,\; B_{3}^{+}=iK\) & 3 & \(+1\) \\
巡回   & \(B_{1}^{-}=I,\; B_{2}^{-}=J,\; B_{3}^{-}=K\)     & 3 & \(-1\) \\
零     & \(N_{0}=e_{4}j,\; N_{1}=e_{4}(kI),\; N_{2}=e_{4}(kJ),\; N_{3}=e_{4}(kK)\) & 4 & \(0\) \\
\hline
\end{tabular}
\end{table}

\section{双対性がローレンツ代数を定める}
\label{sec:duality}

六つの双曲型および巡回生成元は内部で閉じる。同軸のブーストと回転は可換である。隣接する回転は残る回転を再生し、隣接するブースト--回転対は残るブーストを再生し、隣接するブーストは符号を反転して残る回転を再生する：
\begin{align*}
[B^{-}_{a},B^{-}_{a+1}] &= 2B^{-}_{a+2},\\
[B^{+}_{a},B^{-}_{a+1}] &= 2B^{+}_{a+2},\\
[B^{+}_{a},B^{+}_{a+1}] &= -2B^{-}_{a+2}.
\end{align*}
添字は三軸の上を巡回する。これらは回転--ブースト表示 \([J,J]\sim J\)、\([J,K]\sim K\)、\([K,K]\sim -J\) における \(\mathfrak{so}(3,1)\) の構造定数である。最終行の負号こそが、ローレンツ代数をコンパクトな \(\mathfrak{so}(4)\) から分かつ。そこではブースト同士の括弧は \(+J\) となる。

四つの零生成元は、そのローレンツ作用によって互いのあいだで動かされる。動径ブーストは時間並進を自らの軸に沿う空間並進と交換し、残る二つを固定する。回転はその平面内の二つの並進を互いに回し、残りを固定する。代表的な括弧は
\begin{align*}
[B^{+}_{0},N_{0}]&=2N_{1},&
[B^{+}_{0},N_{1}]&=2N_{0},&
[B^{+}_{0},N_{2}]&=0,\\
[B^{-}_{0},N_{2}]&=2N_{3},&
[B^{-}_{0},N_{3}]&=-2N_{2},&
[B^{-}_{0},N_{0}]&=0
\end{align*}
である。そのような括弧はすべて零のままである。零生成元同士は、\eqref{eq:null-strong-vanishing} により可換である。

ブーストと回転の6次元スパンを \(\mathfrak{l}\)、4次元零スパンを \(\mathfrak{n}\) と書く。交換子の双線型性は生成元の恒等式を
\[
[\mathfrak{l},\mathfrak{l}]\subseteq\mathfrak{l},\qquad
[\mathfrak{l},\mathfrak{n}]\subseteq\mathfrak{n},\qquad
[\mathfrak{n},\mathfrak{n}]=0
\]
へ持ち上げる。したがって \(\mathfrak{l}\) と \(\mathfrak{p}:=\mathfrak{l}\oplus\mathfrak{n}\) は \(G(3,1,1)\) のリー部分代数であり、\(\mathfrak{n}\) は \(\mathfrak{p}\) の可換イデアルである。商 \(\mathfrak{p}/\mathfrak{n}\) は \(\mathfrak{l}\) であり、\(\mathfrak{l}\) はいま記録したベクトル表現によって \(\mathfrak{n}\) に作用する。これが、十個のバイベクトルについての代数的事実としての半直積 \(\mathfrak{so}(3,1)\ltimes\mathbb{R}^{3,1}\) である。作用はモーターが要求する意味で忠実である。\([B^{+}_{0},N_{0}]=2N_{1}\neq 0\) であるから、ねじれと並進は一般に可換でない。

\(\mathfrak{l}\) の内部構造は双対性である。写像 \(X\mapsto Xi\) の平方は \(-1\) である。\(i^{2}=-1\) だからである。擬スカラー \(i\) は各ベクトル \(e_{\mu}\) と反可換であり、したがってそのうち二つの任意の積とは可換である。\(i\) は \(\mathfrak{l}\) 上で中心的である。両者を合わせれば、\(i\) による右乗法はローレンツセクター上の複素構造であり、\(\mathfrak{l}\) は \(\mathbb{R}[i]\cong\mathbb{C}\) 上の加群である。

生成元の上では複素構造は厳密である：
\[
B^{-}_{a}\,i \;=\; B^{+}_{a},\qquad
B^{+}_{a}\,i \;=\; -B^{-}_{a}\qquad(a=0,1,2).
\]
すべてのブーストは対応する回転の双対である。\(\mathfrak{r}\) を回転生成元の3次元スパンとすれば、ブーストセクターは \(\mathfrak{r}\,i\) であり、
\[
\mathfrak{l} \;=\; \mathfrak{r}\;\oplus\;\mathfrak{r}\,i
\]
である。ローレンツ代数は回転代数の複素化である。これは古典的な同型 \(\mathfrak{so}(3,1)\cong\mathfrak{so}(3)\otimes\mathbb{C}\cong\mathfrak{sl}(2,\mathbb{C})\) を、\(G(3,1,1)\) の内部で「ブーストは双対化された回転である」という主張として書いたものである。

交換子はこの構造を尊重する。\(i\) が \(\mathfrak{l}\) のすべての元と可換であるから、
\[
[Xi,Y] \;=\; [X,Y]\,i \;=\; [X,Yi],\qquad
[Xi,Yi] \;=\; [X,Y]\,i^{2} \;=\; -[X,Y]
\qquad(X,Y\in\mathfrak{l}).
\]
括弧は複素双線型である。その結果、ローレンツの表の全体が回転の表のみから生成される。\([B^{-}_{a},B^{-}_{a+1}]=2B^{-}_{a+2}\) から
\[
[B^{+}_{a},B^{-}_{a+1}] \;=\; [B^{-}_{a},B^{-}_{a+1}]\,i \;=\; 2B^{+}_{a+2},
\qquad
[B^{+}_{a},B^{+}_{a+1}] \;=\; -[B^{-}_{a},B^{-}_{a+1}] \;=\; -2B^{-}_{a+2}
\]
が得られる。ブースト同士の括弧の負号は \(i^{2}=-1\) である。ユークリッド符号では4次元の擬スカラーは平方が \(+1\) であり、同じ計算は \([K,K]\sim +J\) を与え、複素化はコンパクトな \(\mathfrak{so}(4)\cong\mathfrak{so}(3)\oplus\mathfrak{so}(3)\) へ潰れる。二重時空のローレンツ符号は、双対写像自身に符号化されている。

零セクターは異なる振る舞いをする。\(e_{4}\) が \(i\) と可換であるため、
\[
N_{\mu}\,i \;=\; (e_{4}e_{\mu})\,i \;=\; e_{4}\,(e_{\mu}\,i)
\]
であり、\(e_{\mu}i\) は \(\mathrm{Cl}(3,1)\) において次数3であるから、\(N_{\mu}i\) は \(G(3,1,1)\) において次数4である。双対性は各並進をバイベクトル次数の外へ送る。零イデアルは双対閉ではない。複素構造は \(\mathfrak{p}=\mathfrak{l}\oplus\mathfrak{n}\) のローレンツ因子に属し、ポアンカレスパン全体には属さない。

\section{モーター}
\label{sec:motors}

拡張不整合はねじれバイベクトルと並進バイベクトルの和である、
\begin{equation}
\Omega_{\mathrm{biv}}^{(5)}
= \underbrace{\sum_{a=1}^{3} \Bigl( \frac{\alpha_{a}}{2}\, B_{a}^{+} + \frac{\beta_{a}}{2}\, B_{a}^{-} \Bigr)}_{\Omega_{\mathrm{torsion}}}
\;+\;
\underbrace{\sum_{\mu=0}^{3} \frac{\lambda^{\mu}}{2}\, N_{\mu}}_{\Omega_{\mathrm{trans}}}.
\label{eq:Omega-decomp}
\end{equation}
六つの係数 \(\alpha_{a}\)、\(\beta_{a}\) は通常セクターと双対セクターのもとの不整合である。四つの係数 \(\lambda^{\mu}\) は並進不整合のローレンツ四ベクトルである。光的並進は錐 \(\eta_{\mu\nu}\lambda^{\mu}\lambda^{\nu}=0\) である。

すべての混合積 \(N_{\mu}N_{\nu}\) が消えるため、並進生成元は冪零である：
\[
\Omega_{\mathrm{trans}}^{2}
\;=\;
\frac{1}{4}\sum_{\mu,\nu}\lambda^{\mu}\lambda^{\nu}\,N_{\mu}N_{\nu}
\;=\; 0.
\]
したがってその指数は一次で打ち切られる、
\[
\exp(\Omega_{\mathrm{trans}}) \;=\; 1+\Omega_{\mathrm{trans}}.
\]
\(T(p)=1+\Omega_{\mathrm{trans}}(p)\)、\(T(q)=1+\Omega_{\mathrm{trans}}(q)\) とすれば、積は \(T(p)T(q)=T(p+q)\) である。トランスレータは乗法の下で可換群をなす。並進は純シフトである。曲率も高次項も導入しない。

幾何学的なモーターは二つの指数の順序積である、
\[
M \;:=\; R\,T,
\qquad
R \;:=\; \exp(\Omega_{\mathrm{torsion}}),\qquad
T \;:=\; 1+\Omega_{\mathrm{trans}}.
\]
順序は重要である。交叉交換子 \([B_{a}^{\pm},N_{\mu}]\) は一般に非零であり、したがって \(\Omega_{\mathrm{torsion}}\) と \(\Omega_{\mathrm{trans}}\) は可換である必要がなく、\(M\) は \(\exp(\Omega_{\mathrm{biv}}^{(5)})\) に等しい必要がない。可換なときは二つの式は一致する。可換でないとき---動径ブーストと時間並進が最も単純な例である---幾何学的対象は積 \(RT\) のままである。純動径ブーストはトランスレータをトランスレータへ共役する。サンドイッチ \(R(\varphi)\,T(p)\,\widetilde{R(\varphi)}\) は、係数がブースト平面上のラピディティ \(\varphi\) のローレンツ作用であるトランスレータである。

\(M\) のユニタリ性は次数の帰結である。反転はすべてのバイベクトルをその負へ送り、\(\widetilde{B_{a}^{\pm}}=-B_{a}^{\pm}\) および \(\widetilde{N_{\mu}}=-N_{\mu}\) である。並進を次数1の零ベクトルとして取ったならば、それらは反転偶数となり、下記の論証は失敗する。これが、射影幾何代数が並進を零\emph{バイベクトル}によって生成する理由である~\cite{Gunn2017,Lengyel2024,GunnDeKeninck2020}。ローターに対しては標準的な相殺
\[
R\,\widetilde{R}
\;=\;
\exp(\Omega_{\mathrm{torsion}})\,\exp(-\Omega_{\mathrm{torsion}})
\;=\; 1
\]
が成り立つ。トランスレータに対しては、反転奇数性と冪零性から
\[
T\,\widetilde{T}
\;=\;
\bigl(1+\Omega_{\mathrm{trans}}\bigr)\bigl(1-\Omega_{\mathrm{trans}}\bigr)
\;=\;
1-\Omega_{\mathrm{trans}}^{2}
\;=\; 1
\]
が得られる。反転は反準同型であるから、
\[
M\,\widetilde{M}
\;=\;
R\,T\,\widetilde{T}\,\widetilde{R}
\;=\;
R\,\widetilde{R}
\;=\; 1
\]
である。因子が可換であるか否かに依らず、モーターはユニタリである。

双対性はねじれパラメータの上で四分の一回転として作用する。\(B^{-}_{a}i=B^{+}_{a}\) および \(B^{+}_{a}i=-B^{-}_{a}\) を用いれば、
\[
\Omega_{\mathrm{torsion}}(\alpha,\beta)\,i
\;=\;
\sum_{a}\Bigl(\frac{\beta_{a}}{2}\,B^{+}_{a}-\frac{\alpha_{a}}{2}\,B^{-}_{a}\Bigr)
\;=\;
\Omega_{\mathrm{torsion}}(\beta,-\alpha).
\]
二度施せば、\(i^{2}=-1\) と整合して六つの係数すべてが反転する。\(z_{a}=\alpha_{a}+\mathrm{i}\beta_{a}\) と書けば、双対性は \(-\mathrm{i}\) 倍である。ねじれスカラー
\[
J \;=\; \tfrac12\sum_{a}(\alpha_{a}^{2}-\beta_{a}^{2}) \;=\; \tfrac12\operatorname{Re}\sum_{a}z_{a}^{2}
\]
は正則二次形式の実部であり、符号なしの質量 \(\tfrac12\sum_{a}|z_{a}|^{2}\) はそのエルミート的な相棒である。したがって双対性は \(J\) を反転し、質量を保つ。親理論の通常／双対交換は \(\alpha_{a}\) と \(\beta_{a}\) を入れ替え、この四分の一回転とは新しい回転パラメータの符号だけが異なる。両者はともに \(J\) を反転する。並進パラメータは別に追跡しなければならない。双対性が各 \(N_{\mu}\) を次数4へ送るからである。

10次元空間全体の上の二次診断量は、\(\mathfrak{iso}(3,1)\) のカルタン--キリング形式だけでは得られない。その形式は並進の上で消える~\cite{FultonHarris1991,Humphreys1972}。ローレンツセクターのキリング二次形式を並進係数上のミンコフスキー計量によって拡張すると、非退化な対
\begin{equation}
J^{(5)}
\;:=\;
\tfrac{1}{16}\,B_{\mathrm{Killing}}(\Omega_{\mathrm{torsion}},\Omega_{\mathrm{torsion}})
\;+\;
\tfrac{1}{2}\,\eta_{\mu\nu}\lambda^{\mu}\lambda^{\nu}
\label{eq:J5-invariant}
\end{equation}
が得られる。\eqref{eq:Omega-decomp} の半角展開のもとで、\(B(B_{a}^{+},B_{a}^{+})=+8\) および \(B(B_{a}^{-},B_{a}^{-})=-8\) とすると、第一項は \(\tfrac18\sum_{a}(\alpha_{a}^{2}-\beta_{a}^{2})\) に等しく、これは離散的な伴う論文の \(J=\tfrac12\sum_{a}(\alpha_{a}^{2}-\beta_{a}^{2})\) の四分の一である。第二項は光錐の上でまさに消える。二つの寄与が分離されるのは、\(\mathfrak{iso}(3,1)\) のキリング形式が半単純生成元と並進生成元のあいだで消えるからである。並進寄与は \(\lambda^{\mu}\) が制約されない限り非有界である。離散伴論文の許容天井 \(|J|\le 3\pi^{2}/8\) はねじれ部分にだけ付く。

モーターの算法には必要ないが、同じ代数の幾何学的な読みを一つ記録しておく。\(G(3,1,1)\) において平面は三ブレードである。それらの幾何積の反対称部分
\[
A \mathbin{\underline{\vee}} B \;:=\; \tfrac12(AB-BA)
\]
は次数2に値をとる。代表的な構成の上で、\(P\) が三ブレードであり、\(D\) が \(P\) の平面の法線 \(n\) に沿う双対セクターの要素であるとき、括弧積はその法線に沿う零バイベクトルである、
\[
P \mathbin{\underline{\vee}} D \;=\; \kappa\,(e_{4}\wedge n).
\]
双対時空は法線の担体として作用し、括弧積は法線に沿う並進を生じる。付録はそのような例を一つ計算する。これらの括弧積を \(\Omega_{\mathrm{biv}}^{(5)}\) に結ぶ完全に共変な算法は、ここでは展開しない。

\section{チャート上の重力}
\label{sec:gravity}

親理論は遠平行重力の実現を目指す。幾何は曲率テンソルではなくねじれによって担われる。射影拡張において、同じ考えは固定座標チャートの上で利用できる。終始第\ref{sec:duality}節のポアンカレ部分代数を扱う。一般多様体上の遠平行スカラーとアインシュタイン--ヒルベルト作用との変分等価性は TEGR 文献~\cite{maluf2013} から取り、再導出しない。

ユニタリモーターは、サンドイッチによって参照コフレームに作用する。\(G(3,1,1)\) 内部の \(\mathrm{Cl}(3,1)\) ベクトル生成元 \(\{\mathring{e}^{a}\}_{a=0}^{3}\) を固定し、
\begin{equation}
e^{a} \;:=\; \bigl\langle M\,\mathring{e}^{a}\,\widetilde{M}\bigr\rangle_{1}
\label{eq:motor-tetrad}
\end{equation}
とおく。チャート \((x^{\mu})\) 上で成分場 \(e^{a}{}_{\mu}(x)\) は計量
\begin{equation}
g_{\mu\nu} \;=\; \eta_{ab}\, e^{a}{}_{\mu}\, e^{b}{}_{\nu}
\label{eq:induced-metric}
\end{equation}
を誘導する。ここで \(\eta=\mathrm{diag}(-1,+1,+1,+1)\) である。\(M\) が純ローレンツローターであるとき、サンドイッチは参照ベクトル上の局所ローレンツ作用である。零並進は完全モーター \(M=RT\) に入るが、それ自体では古典的シュヴァルツシルト対角ゲージを供給しない。それらは \(J^{(5)}\) によって測られる並進不整合として利用可能なままである。

チャート上のコフレームのヴァイツェンベックねじれは反対称化された導関数
\begin{equation}
T^{a}{}_{\mu\nu} \;=\; \partial_{\mu} e^{a}{}_{\nu} - \partial_{\nu} e^{a}{}_{\mu}
\label{eq:weitzenbock-torsion}
\end{equation}
である。逆テトラッド \(e_{a}{}^{\mu}\) および \(g_{\mu\nu}\) によって下げられた時空添字のもとで、遠平行スカラーは
\begin{equation}
T \;=\; \tfrac14 T_{\rho\mu\nu}T^{\rho\mu\nu} + \tfrac12 T_{\rho\mu\nu}T^{\nu\mu\rho} - T_{\rho}T^{\rho}
\label{eq:teleparallel-T}
\end{equation}
である。ここで \(T_{\rho}=T^{\sigma}{}_{\rho\sigma}\) である。計量適合 Levi-Civita 接続の曲率テンソルは不要である。遠平行記述は幾何を \(T^{a}{}_{\mu\nu}\) に格納する。

外部シュヴァルツシルト幾何は、この構成の動径ブーストである。\(r_{s}=2GM/c^{2}>0\) とし、領域 \(r>r_{s}\) で作業する。対角コフレーム（単位 \(c=1\)）は
\begin{align}
e^{0} &= \sqrt{1-\frac{r_{s}}{r}}\,dt,\nonumber\\
e^{1} &= \bigl(1-\frac{r_{s}}{r}\bigr)^{-1/2}\,dr,\nonumber\\
e^{2} &= r\,d\theta,\nonumber\\
e^{3} &= r\sin\theta\,d\phi
\label{eq:schwarzschild-tetrad}
\end{align}
である。\(A(r):=1-r_{s}/r\) とおき、
\begin{equation}
\varphi(r) \;:=\; -\log\sqrt{A(r)} \;=\; -\tfrac12\log A(r)
\qquad(r>r_{s})
\label{eq:radial-rapidity}
\end{equation}
とする。このとき \(e^{-\varphi}=\sqrt{A}\) および \(e^{\varphi}=A^{-1/2}\) であり、したがって \((t,r)\) 対角係数は \(B^{+}_{1}=e_{0}\wedge e_{1}\) によって生成される純双曲型ローターのブーストスケール因子である。角脚は球面参照コフレームから来る。対応するねじれパラメータは純ブースト \(\Omega_{\mathrm{torsion}}=(\varphi/2)B^{+}_{1}\) であり、対角ゲージでは零パラメータは消える。誘導計量~\eqref{eq:induced-metric} はシュヴァルツシルト計量である。

いま同じチャートの上に三つのスカラーがあり、それらは同じ対象ではない。有限角の診断量
\[
J(\varphi) \;=\; \tfrac12\varphi(r)^{2}
\]
はパラメータ空間におけるブーストの大きさを測る。場の種
\begin{equation}
J_{\mathrm{field}}(r) \;:=\; \tfrac12\bigl(\varphi'(r)\bigr)^{2},
\label{eq:J-field}
\end{equation}
ただし \(\varphi'=-r_{s}/(2r^{2}A)\)、はそのブーストの動径変化率を測る。遠平行スカラー~\eqref{eq:teleparallel-T} はコフレームの導関数の二次形式である。この外部チャート上では
\[
T \;=\; \frac{2}{r^{2}}\frac{(1-\sqrt{A})^{2}}{\sqrt{A}}
\]
と評価され、TEGR 動径カレント \(B^{r}=4(1-\sqrt{A})/r\) のもとで純動径発散
\[
T \;=\; -r^{-2}\,\partial_{r}(r^{2}B^{r})
\]
である。恒等式 \(\partial_{r}\sqrt{A}=-\varphi'\sqrt{A}\) はブーストねじれ脚を \(\varphi'\) に結ぶが、\(J\) や \(J_{\mathrm{field}}\) を \(T\) と同一視はしない。

定数の非零ブーストは \(J\neq 0\) をもち、平坦チャートは \(T=0\) をもつ。ミンコフスキー空間上の純空間並進は \(J^{(5)}\neq 0\) かつ \(T=0\) である。場の種でさえ、一般の外部点で \(\tfrac12 T\) と一致しない。たとえば \((r_{s},r)=(2,4)\) において二つの数は食い違う。有限角、角の変化率、コフレーム密度は、導関数の梯子の異なる段に住む。一般の滑らかなモーター場 \(M(x)=R(x)T(x)\) に対する辞書は、有限角ではなく \((\Omega_{\mathrm{torsion}},\lambda)\) の導関数から組まれなければならず、全発散を含むことが期待される。その辞書はここでは構成しない。

\(G(3,1,1)\) の完全な退化計量に関するサンドイッチも同様に脇に置く。\(M\) と \(\exp(\Omega_{\mathrm{biv}}^{(5)})\) との同一視は因子が可換なとき成り立ち、そうでないときは一般に成り立たない。

\section{結論}
\label{sec:conclusions}

二重時空の双四元数代数に零ベクトル \(e_{4}\) を付加すると、並進はバイベクトルとして実現される。\(G(3,1,1)\) の10次元バイベクトル空間はポアンカレ代数 \(\mathfrak{so}(3,1)\ltimes\mathbb{R}^{3,1}\) である。六つのねじれ生成元は閉じ、四つの零生成元はそれらが作用する可換イデアルを張る。双対性 \(X\mapsto Xi\) はローレンツ因子の複素構造である。すべてのブーストは双対化された回転であり、括弧は複素双線型であり、\(\mathfrak{so}(3,1)\) を \(\mathfrak{so}(4)\) から分ける符号は \(i^{2}=-1\) である。並進は次数4へ送られる。双対性は零イデアルを閉じない。

両セクターを一つにするモーターは順序積 \(M=RT\) である。零生成元の冪零性はその指数をトランスレータ \(T=1+\Omega_{\mathrm{trans}}\) へ打ち切る。すべてのバイベクトルの反転奇数性は \(M\widetilde{M}=1\) を与える。ねじれと並進が可換でないとき、積は単一の指数ではない。それは幾何学的対象のままである。双対性はねじれパラメータを四分の一回転し、\(J\) を反転し、符号なしの質量を保つ。診断量 \(J^{(5)}\) はキリング二次形式を並進上のミンコフスキー対によって拡張する。その並進項は光錐の外で非有界である。

チャート上では、\(M\) によるサンドイッチがテトラッドを生じる。外部シュヴァルツシルトコフレームはそのテトラッドの動径ブーストである。有限角のスカラー \(J\)、場の種 \(J_{\mathrm{field}}\)、遠平行密度 \(T\) は同一幾何の三つの異なる読みであり、点ごとに同一ではない。残るのは、平面に対する共変な括弧積算法、\(G(3,1,1)\) の退化計量によるサンドイッチ、および一般モーターに対する場の辞書である。TEGR とアインシュタイン--ヒルベルト重力の変分等価性は~\cite{maluf2013} の入力である。

\appendix

\section{並進を生じる平面括弧積}
\label{app:bracket}

\(G(3,1,1)\) において平面は三ブレードである。二つの三ブレード \(A\) および \(B\) に対し、幾何積の反対称部分および対称部分は
\[
A \mathbin{\underline{\vee}} B \;=\; \tfrac12(AB-BA),\qquad
A \mathbin{\overline{\vee}} B \;=\; \tfrac12(AB+BA)
\]
である。これらの記号はこの意味でのみ用いる。PGA 文献の反内積および反外積ではない~\cite{Lengyel2024,DorstFontijneMann2007}。5次元ベクトル空間において、二つの次数3要素の幾何積は次数 \(0\)、\(2\)、\(4\) のみをもつ。

\(P=e_{4}\wedge e_{1}\wedge e_{2}\) および親理論の双対セクター三ベクトル \(D=jI=-e_{0}e_{3}e_{2}\) を取る。積 \(PD-DP\) を通じて生成元を反可換に動かすと、括弧積は \(e_{4}\wedge e_{3}\) の倍数に還元される。方向 \(e_{3}\) は空間平面 \(e_{1}\wedge e_{2}\) への単位法線であるから、実係数 \(\kappa\) に対して
\[
P \mathbin{\underline{\vee}} D \;=\; \kappa\, N_{3}
\]
である。右辺はその法線に沿う並進の生成元である。双対要素 \(D\) は法線方向の担体として作用した。

第二の、純空間的な例は次数パターンを確認する。\(A=e_{1}e_{2}e_{3}\) および \(B=e_{1}e_{2}e_{4}\) に対して \(AB=-e_{3}e_{4}\)、\(BA=+e_{3}e_{4}\) となり、したがって \(A\mathbin{\underline{\vee}}B=-e_{3}e_{4}\)、\(A\mathbin{\overline{\vee}}B=0\) である。反対称積に期待されるとおり、純次数2の交換子である。

\section{Lean で検証した恒等式}
\label{app:lean}

{\raggedright
\emergencystretch=3em

本稿の構成は Lean~4 プロジェクト \texttt{DstDiophantine} において形式化されている。以下は第二の証明ではない。導入の三つの主張を、それを検査する定理へ結ぶ地図である。証明は指名したモジュールにあり、ここでは再現しない。十個の生成元の線形独立性、\(\mathfrak{iso}(3,1)\) との抽象的な同型、一般モーターに対する場の辞書、および \(G(3,1,1)\) の完全な退化計量によるサンドイッチは、形式化では主張していない。

\subsection{零方向と十個の生成元}

代数は、重み \((-1,+1,+1,+1,0)\) の二次形式に対するクリフォード代数であり、モジュールは \texttt{Algebra.QuadraticForm} および \texttt{Algebra.PGA} である。恒等式 \(e_{4}^{2}=0\) および \(\{e_{4},e_{\mu}\}=0\) は \texttt{PGA.e4\_sq} と \texttt{PGA.e4\_anticomm} である。表~\ref{tab:generators} の十個のバイベクトルは \texttt{Algebra.Generators} の \texttt{hyperbolic}、\texttt{cyclic}、\texttt{null} である。平方と強い消滅は
{\small
\begin{verbatim}
theorem hyperbolic_sq (a : Fin 3) : hyperbolic a * hyperbolic a = 1
theorem cyclic_sq (a : Fin 3) : cyclic a * cyclic a = -1
theorem null_mul_null (mu nu : Fin 4) : null mu * null nu = 0
theorem null_reverse (mu : Fin 4) : reverse (null mu) = -null mu
\end{verbatim}
}
最後の行が、並進がバイベクトルでなければならない理由である。次数1の零ベクトルの反転は偶数となり、第\ref{sec:motors}節のユニタリ性の論証は失敗する。同じモジュールが \texttt{hyperbolic\_reverse} および \texttt{cyclic\_reverse} を証明する。

\subsection{双対性とポアンカレスパン}

\texttt{Algebra.LorentzLie} は生成元の表をスパンへ持ち上げる。\(\mathfrak{l}\) を \texttt{lorentzSpan}、\(\mathfrak{n}\) を \texttt{nullSpan} と書く。ローレンツスパンと10次元スパンは交換子の下で閉じる（\texttt{commutator\_mem\_lorentzSpan}、\texttt{commutator\_mem\_poincareSpan}）。零スパンは可換イデアルである（\texttt{commutator\_null\_null\_eq\_zero}、\texttt{commutator\_poincare\_null\_mem\_nullSpan}）。これらは mathlib のリー部分代数 \texttt{lorentzLieSubalgebra} および \texttt{poincareLieSubalgebra} として与えられる。作用はモーターが要求する意味で忠実である。\texttt{commutator\_hyperbolic0\_null0\_ne\_zero} は \([B^{+}_{0},N_{0}]\neq 0\) である。

双対性は Cl(3,1) 擬スカラーによる右乗法である。ローレンツスパンの上では複素構造であり、括弧は複素双線型である：
{\small
\begin{verbatim}
theorem hyperbolic_eq_dual_cyclic (a : Fin 3) :
    hyperbolic a = dual (cyclic a)
theorem dual_dual (x : PGA) : dual (dual x) = -x
theorem commutator_dual_dual (x y : PGA)
    (hx : x Mem lorentzSpan) (hy : y Mem lorentzSpan) :
    commutator (dual x) (dual y) = -commutator x y
theorem commutator_hyperbolic_hyperbolic_cyclic (a b : Fin 3) :
    commutator (hyperbolic a) (hyperbolic b) =
      if a = b then 0
      else if b = a + 1 then -((2 : Real) smul cyclic (a + 2))
      else (2 : Real) smul cyclic (b + 2)
\end{verbatim}
}
ブースト同士の表の負号は \(i^{2}=-1\) である。すべてのローレンツ要素は回転と双対化された回転の和である（\texttt{mem\_lorentzSpan\_iff}）。ねじれパラメータの上で双対性は四分の一回転 \((\alpha,\beta)\mapsto(\beta,-\alpha)\) である（\texttt{dual\_omegaTorsion}）。それは \(J\) を反転し、符号なしの質量を保つ（\texttt{J\_dualParams}、\texttt{mass\_dualParams}）。

\subsection{モーター}

\texttt{Algebra.Motor} は \(\Omega_{\mathrm{torsion}}\)、\(\Omega_{\mathrm{trans}}\)、および順序積 \(RT\) を定義する。並進生成元の冪零性は Banach 指数を打ち切り、トランスレータは加法し、積はユニタリである：
{\small
\begin{verbatim}
theorem omegaTrans_sq (p : TransParams) :
    omegaTrans p * omegaTrans p = 0
theorem exp_omegaTrans (p : TransParams) :
    exp (omegaTrans p) = 1 + omegaTrans p
theorem expTrans_mul (p q : TransParams) :
    expTrans p * expTrans q =
      expTrans (fun mu => p.lambda mu + q.lambda mu)
theorem rotor_unitary (p : TorsionParams) :
    rotorTorsion p * reverse (rotorTorsion p) = 1
theorem motor_unitary (p : OmegaParams) :
    motor p * reverse (motor p) = 1
\end{verbatim}
}
因子が可換なとき \(\exp(\Omega_{\mathrm{biv}}^{(5)})=RT\) である（\texttt{exp\_omegaBiv\_eq\_motor\_of\_commute}）。可換でないとき両者は異なる。\texttt{exists\_omegaBiv\_ne\_motor} は動径ブーストと時間並進を示す。純動径ブーストに限らず、すべてのねじれローターがトランスレータをトランスレータへ共役する（\texttt{Algebra.MotorGroup}、\texttt{exists\_sandwich\_rotorTorsion\_expTrans}）。閉じた形 \(\exp(tB^{+})=\cosh t+\sinh t\cdot B^{+}\) および \(\exp(tB^{-})=\cos t+\sin t\cdot B^{-}\) は \texttt{exp\_of\_sq\_one} および \texttt{exp\_of\_sq\_neg\_one} である。

第\ref{sec:motors}節の診断量は \texttt{Invariant.J5} である。伴う論文のスカラー \(J=\tfrac12\sum_{a}(\alpha_{a}^{2}-\beta_{a}^{2})\) と \(\lambda^{\mu}\) 上のミンコフスキー対との和である。並進に制約がなければ非有界である（\texttt{J5\_unbounded}）。

\subsection{チャート上の重力}

対角シュヴァルツシルトコフレームはシュヴァルツシルト計量を誘導する（\texttt{Gravity.Schwarzschild}、\texttt{inducedMetric\_schwarzschild}）。\((t,r)\) 脚は動径ラピディティ \(\varphi=-\tfrac12\log A\) のブーストスケール因子であり（\texttt{Gravity.Sandwich}、\texttt{boostScaleFactors\_schwarzschild}）、その切片上の有限角診断量は \(J=\tfrac12\varphi^{2}\) である（\texttt{J\_radialBoostParams}）。

\texttt{Gravity.JTDictionary} は同一チャート上のヴァイツェンベックスカラーを
\[
T \;=\; \frac{4}{r^{2}}(\cosh\varphi-1)
\]
へ変換する。これは第\ref{sec:gravity}節の \(\frac{2}{r^{2}}(1-\sqrt{A})^{2}/\sqrt{A}\) と同一である（\texttt{schwarzschild\_T\_eq\_cosh\_rapidity}）。有限角、場の種、コフレーム密度は別物のままである。その節の三つの比較は \texttt{Gravity.Identification} において検査される。非零の有限角と平坦チャート、純空間並進と消滅する \(T\)、および \((r_{s},r)=(2,4)\) における場の種と \(\tfrac12 T\) である。

\par}

\begin{thebibliography}{99}

\bibitem{Gunn2017}
C.~Gunn,
\textit{Geometric algebras for Euclidean geometry},
Advances in Applied Clifford Algebras \textbf{27} (1), 185--208 (2017).

\bibitem{Lengyel2024}
E.~Lengyel,
\textit{Projective Geometric Algebra Illuminated},
Terathon Software LLC (2024).

\bibitem{DeKeninckDorst2019}
S.~De Keninck and L.~Dorst,
\textit{Projective geometric algebra: A new framework for doing Euclidean geometry},
arXiv:1901.05873 (2019).

\bibitem{GunnDeKeninck2020}
C.~Gunn and S.~De Keninck,
\textit{Geometric algebra and computer graphics},
in: SIGGRAPH 2019 Courses, ACM, New York (2019).
See also: \texttt{https://bivector.net/}.

\bibitem{DorstFontijneMann2007}
L.~Dorst, D.~Fontijne, and S.~Mann,
\textit{Geometric Algebra for Computer Science: An Object-Oriented Approach to Geometry},
Morgan Kaufmann, San Francisco (2007).

\bibitem{Hestenes2015}
D.~Hestenes,
\textit{Space-Time Algebra}, 2nd ed.,
Birkh\"{a}user, Basel (2015).

\bibitem{HestenesSobczyk1984}
D.~Hestenes and G.~Sobczyk,
\textit{Clifford Algebra to Geometric Calculus: A Unified Language for Mathematics and Physics},
Reidel, Dordrecht (1984).

\bibitem{Lounesto2001}
P.~Lounesto,
\textit{Clifford Algebras and Spinors}, 2nd ed.,
London Mathematical Society Lecture Note Series \textbf{286},
Cambridge University Press, Cambridge (2001).

\bibitem{ConwaySmith2003}
J.~H.~Conway and D.~A.~Smith,
\textit{On Quaternions and Octonions: Their Geometry, Arithmetic, and Symmetry},
A K Peters, Natick, MA (2003).

\bibitem{DoranLasenby2003}
C.~Doran and A.~Lasenby,
\textit{Geometric Algebra for Physicists},
Cambridge University Press, Cambridge (2003).

\bibitem{Selig2005}
J.~M.~Selig,
\textit{Geometric Fundamentals of Robotics}, 2nd ed.,
Monographs in Computer Science, Springer, New York (2005).

\bibitem{Weinberg1995}
S.~Weinberg,
\textit{The Quantum Theory of Fields, Volume I: Foundations},
Cambridge University Press, Cambridge (1995).

\bibitem{Tung1985}
W.-K.~Tung,
\textit{Group Theory in Physics},
World Scientific, Singapore (1985).

\bibitem{FultonHarris1991}
W.~Fulton and J.~Harris,
\textit{Representation Theory: A First Course},
Graduate Texts in Mathematics \textbf{129},
Springer, New York (1991).

\bibitem{Humphreys1972}
J.~M.~Humphreys,
\textit{Introduction to Lie Algebras and Representation Theory},
Graduate Texts in Mathematics \textbf{9},
Springer, New York (1972).

\bibitem{maluf2013}
J.~W.~Maluf,
\textit{The teleparallel equivalent of general relativity},
Annalen der Physik \textbf{525} (5--6), 339--357 (2013).

\end{thebibliography}

\end{document}
